\subsection{单实变量函数}
\label{sec:chap2-one-real-variable}
\renewcommand{\thestatement}{4.\number\numexpr\value{statement}-58\relax}

本节主要回顾弱导数概念与通常意义下的导数概念之间的联系. 对于我们关心的单实变量函数, 我们将以有限但充分的方式来完成这一回顾.

本节最后回顾 Gronwall 型不等式, 它们是获得演化偏微分方程解的先验估计的有用工具.

\subsubsection{微分与原函数}
\label{subsec:chap2-differentiation-antiderivatives}

设 $[a,b]$ 是 $\mathbb{R}$ 的一个紧区间. 回顾 $W^{1,1}((a,b))$ 是 $L^1((a,b))$ 中分布意义下的导数仍为 $L^1((a,b))$ 函数的那些函数的集合 (关于 Sobolev 空间的更完整研究见 Chapter~\ref{chap:sobolev-spaces}). 对这样的函数, 可以提出一个基本问题: 它是否在通常意义下可微, 以及能否写出微积分基本定理
\[
 f(y)=f(x)+\int_x^y f'(t)\,\mathrm{d}t,
 \qquad\forall x,y\in[a,b].
\]

这里回顾与该问题有关的一些结果. 在本书后文, 为了证明 Navier--Stokes 方程弱解的动能随时间演化关系成立, 这些材料十分有用 (特别参见\MBUnresolvedReference{sec:chap5-weak-solution-kinetic-energy}{Section 1.4 of Chapter V}).

\begin{Lemma}
\label{lem:chap2-integral-antiderivative}
设 $g\in L^1((a,b))$, $C\in\mathbb{R}$, 并定义
\[
 f(t)=C+\int_a^t g(s)\,\mathrm{d}s.
\]
则 $f$ 在 $[a,b]$ 上连续, $f\in W^{1,1}((a,b))$, 且其分布导数为 $g$.
\end{Lemma}

\begin{Proof}
对 $t_0\in[a,b)$ 和充分小的 $h>0$,
\[
 f(t_0+h)-f(t_0)
 =\int_{t_0}^{t_0+h}g(s)\,\mathrm{d}s
 =\int_a^b g(s)\mathbf{1}_{[t_0,t_0+h]}(s)\,\mathrm{d}s
 \longrightarrow0,
\]
这是 Lebesgue 控制收敛定理的结果. 对 $h<0$ 使用类似论证, 可得 $f$ 的连续性.

现在需要验证 $f$ 的分布意义下的导数是函数 $g$. 设 $\varphi\in\mathcal{D}((a,b))$; 有
\[
 -\int_a^b f(t)\varphi'(t)\,\mathrm{d}t
 =-C\int_a^b\varphi'(t)\,\mathrm{d}t
 -\int_a^b\left(\int_a^t g(s)\varphi'(t)\,\mathrm{d}s\right)\mathrm{d}t.
\]
第一项为零, 因为 $\varphi(a)=\varphi(b)=0$. 对第二项应用 Fubini 定理 (函数 $(t,s)\mapsto\varphi'(t)g(s)$ 关于两个变量可积), 得
\begin{align*}
 -\int_a^b f(t)\varphi'(t)\,\mathrm{d}t
 &=-\int_a^b\int_a^b\mathbf{1}_{[a\leq s\leq t]}g(s)\varphi'(t)\,\mathrm{d}s\,\mathrm{d}t\\
 &=-\int_a^b\int_a^b\mathbf{1}_{[a\leq s\leq t]}g(s)\varphi'(t)\,\mathrm{d}t\,\mathrm{d}s\\
 &=-\int_a^b g(s)\int_s^b\varphi'(t)\,\mathrm{d}t\,\mathrm{d}s\\
 &=-\int_a^b g(s)\bigl(\varphi(b)-\varphi(s)\bigr)\,\mathrm{d}s\\
 &=\int_a^b g(s)\varphi(s)\,\mathrm{d}s.
\end{align*}
这确实证明 $g=f'$ 在分布意义下成立, 因而 $f\in W^{1,1}((a,b))$.
\end{Proof}

\begin{Corollary}
\label{cor:chap2-w11-continuous-representative}
每个 $f\in W^{1,1}((a,b))$ 都几乎处处等于 $[a,b]$ 上某个连续函数 $\widetilde f$, 且
\[
 \widetilde f(y)=\widetilde f(x)+\int_x^y f'(s)\,\mathrm{d}s,
 \qquad\forall x,y\in[a,b].
\]
换言之, 对几乎处处的 $x,y\in[a,b]$, 有
\[
 f(y)=f(x)+\int_x^y f'(s)\,\mathrm{d}s.
\]
\end{Corollary}

\begin{Proof}
设 $f\in W^{1,1}((a,b))$, 并引入
\[
 g(t)=\int_a^t f'(s)\,\mathrm{d}s.
\]
\begin{itemize}
\item 根据 Lemma~\ref{lem:chap2-integral-antiderivative}, 函数 $g$ 连续, 属于 $W^{1,1}((a,b))$, 且其分布意义下的导数是 $f'$. 因而 $f-g$ 的分布导数为零. 我们知道, 存在实数 $C$, 使 $f-g=C$ 几乎处处成立 (这一结果的更一般情形见 Lemma~\ref{lem:chap2-zero-gradient-distribution-constant}). 若定义 $\widetilde f=C+g$, 就证明了 $f$ 几乎处处与连续函数 $\widetilde f$ 相等.
\item 由 $\widetilde f$ 的定义, 显然对所有 $x,y\in[a,b]$ 有
\begin{align*}
 \widetilde f(y)-\widetilde f(x)
 &=\left(C+\int_a^y f'(s)\,\mathrm{d}s\right)
 -\left(C+\int_a^x f'(s)\,\mathrm{d}s\right)\\
 &=\int_x^y f'(s)\,\mathrm{d}s.
\end{align*}
\end{itemize}
\end{Proof}

对 $t_0\in[a,b]$, 以 $V_\eta(t_0)$ 表示 $[a,b]$ 中包含 $t_0$ 且测度小于 $\eta$ 的开邻域全体.

\begin{Definition}[Lebesgue points]
\label{def:chap2-lebesgue-point}
设 $f\in L^1((a,b))$, $t_0\in(a,b)$. 若
\[
 \sup_{\omega\in V_\eta(t_0)}\frac1{|\omega|}
 \int_\omega|f(t)-f(t_0)|\,\mathrm{d}t\longrightarrow0
 \quad(\eta\to0),
\]
则称 $t_0$ 为 $f$ 的 Lebesgue 点.
\end{Definition}

有了这个定义, 可得如下结果.

\begin{Proposition}
\label{prop:chap2-antiderivative-at-lebesgue-point}
设 $f\in L^1((a,b))$, 且 $t_0$ 是 $f$ 的 Lebesgue 点. 则任意原函数
\[
 F(t)=C+\int_a^t f(s)\,\mathrm{d}s
\]
在 $t_0$ 按经典意义可微, 且 $F'(t_0)=f(t_0)$.
\end{Proposition}

\begin{Proof}
只需写出
\begin{align*}
 \left|\frac{F(t_0+h)-F(t_0)}h-f(x_0)\right|
 &=\left|\frac1h\int_{t_0}^{t_0+h}\bigl(f(s)-f(t_0)\bigr)\,\mathrm{d}s\right|\\
 &\leq\frac1h\left|\int_{t_0}^{t_0+h}|f(s)-f(t_0)|\,\mathrm{d}s\right|;
\end{align*}
由 Lebesgue 点的定义, 最后一项在 $h$ 趋于 $0$ 时趋于 $0$.
\end{Proof}

本节的基本定理是测度论中一个相当困难的结果, 这里不予证明. 不过, 例如可在\MBUnresolvedReference{bib:lebesgue-point-sources}{[69, 100]}中找到证明.

\begin{Theorem}
\label{thm:chap2-almost-every-point-lebesgue}
若 $f\in L^1((a,b))$, 则 $(a,b)$ 中几乎每一点都是 $f$ 的 Lebesgue 点.
\end{Theorem}

前面两个结果的直接推论是, $L^1((a,b))$ 中任意函数 $f$ 的原函数 $F$ 几乎处处可微, 且满足 $F'=f$ 以及微积分基本定理
\[
 F(y)=F(x)+\int_x^y F'(t)\,\mathrm{d}t
 =F(x)+\int_x^y f(t)\,\mathrm{d}t,
 \qquad\forall x,y\in[a,b].
\]

最后, 下面这个初等结果十分有用.

\begin{Proposition}
\label{prop:chap2-continuity-point-lebesgue}
设 $f\in L^1((a,b))$. $f$ 的每个连续点都是 Lebesgue 点. 特别地, $f$ 的任意原函数在每个连续点 $t_0$ 可微, 且导数为 $f(t_0)$.
\end{Proposition}

\par\noindent\textbf{Remark 4.1:}\quad 可以证明\MBUnresolvedReference{bib:source-69}{[69]}, $W^{1,1}((a,b))$ 中的函数恰好是 $(a,b)$ 上的绝对连续函数.

以如下结果及其推论结束本小节.

\begin{Lemma}[Hardy's inequality]
\label{lem:chap2-hardy-one-dimensional}
对 $1<p<+\infty$、非负 $f\in L^p((0,+\infty))$ 以及任意 $M\in[0,+\infty]$,
\[
 \int_0^M\left(\frac1x\int_0^x f(s)\,\mathrm{d}s\right)^p\mathrm{d}x
 \leq\left(\frac{p}{p-1}\right)^p\int_0^M f(s)^p\,\mathrm{d}s.
\]
\end{Lemma}

注意, 当 $p=1$ 时, 类似的不等式不成立.

\begin{Proof}
先证明 $M=+\infty$ 时的不等式. 一般情形可通过取 $f_M(s)=\mathbf{1}_{[0,M]}(s)f(s)$ 得到.

由稠密性, 只需对 $f\in C_c^\infty((0,+\infty))$ 证明该结果. 对这样的 $f$, 令
\[
 F(x)=\frac1x\int_0^x f(s)\,\mathrm{d}s,
\]
并注意到 $F$ 在 $0$ 的某个邻域中为零, 且当 $x$ 充分大时, 对某个 $C\in\mathbb{R}$ 有 $F(x)=C/x$. 特别地, $F\in L^p((0,+\infty))$. 注意
\[
 \frac{\mathrm{d}(xF(x))}{\mathrm{d}x}=f(x),
\]
因而可如下分部积分:
\begin{align*}
 \int_0^{+\infty}F(x)^p\,\mathrm{d}x
 &=\int_0^{+\infty}(xF(x))^p\frac1{x^p}\,\mathrm{d}x\\
 &=\frac{p}{p-1}\int_0^{+\infty}(xF(x))^{p-1}f(x)\frac1{x^{p-1}}\,\mathrm{d}x\\
 &=\frac{p}{p-1}\int_0^{+\infty}F(x)^{p-1}f(x)\,\mathrm{d}x.
\end{align*}
最后使用 Hölder 不等式:
\[
 \int_0^{+\infty}F(x)^p\,\mathrm{d}x
 \leq\frac{p}{p-1}\|F\|_{L^p}^{p-1}\|f\|_{L^p}.
\]
\end{Proof}

由这一不等式可立即推出下面的结果, 它在后文中十分重要 (特别参见 Proposition~\ref{prop:chap3-hardy-boundary-distance} 的多维版本).

\begin{Corollary}
\label{cor:chap2-hardy-boundary-quotient}
设 $f\in W^{1,p}((a,b))\subset W^{1,1}((a,b))$, $1<p<+\infty$, 且 $f(a)=0$. 则 $g(x)=f(x)/(x-a)$ 属于 $L^p((a,b))$, 且
\[
 \|g\|_{L^p}\leq C\|f'\|_{L^p}.
\]
\end{Corollary}

\subsubsection{微分不等式与 Gronwall 引理}
\label{subsec:chap2-gronwall}

下面这些关于常微分方程和不等式的引理, 是研究含时偏微分方程, 特别是证明能量估计时非常有用的工具.

\begin{Lemma}
\label{lem:chap2-linear-decay}
设 $\alpha>0$, $\beta\geq0$, $y\in C^1([0,+\infty),\mathbb{R})$ 且
\[
 y'(t)+\alpha y(t)\leq\beta,
 \qquad\forall t\geq0.
\]
则
\[
 y(t)\leq y(0)e^{-\alpha t}+\frac\beta\alpha,
 \qquad\forall t\geq0.
\]
\end{Lemma}

\begin{Proof}
将不等式乘以 $e^{\alpha t}$ 后积分.
\end{Proof}

下面这个引理称为 Gronwall 不等式, 尽管其当前形式的证明归功于 Bellman\MBUnresolvedReference{bib:source-15}{[15]}. 它在证明 (偏) 微分方程解的先验估计时处于核心地位.

\begin{Lemma}
\label{lem:chap2-gronwall}
设 $y\in L^\infty((0,T))$, $g\in L^1((0,T))$ 非负, $y_0\in\mathbb{R}$, 且对几乎处处的 $t\in(0,T)$,
\[
 y(t)\leq y_0+\int_0^t g(s)y(s)\,\mathrm{d}s.
\]
则
\[
 y(t)\leq y_0\exp\left(\int_0^t g(s)\,\mathrm{d}s\right)
\]
对几乎处处的 $t\in(0,T)$ 成立.
\end{Lemma}

\begin{Proof}
令
\[
 h(t)=y_0+\int_0^t g(s)y(s)\,\mathrm{d}s,
\]
即假设中不等式的右端. 由于 $y\in L^\infty((0,T))$ 且 $g\in L^1((0,T))$, 函数 $h$ 属于 $W^{1,1}((0,T))$, 因而几乎处处可微, 且其导数为 $gy$ (见 Section~\ref{subsec:chap2-differentiation-antiderivatives}). 此外, 由假设以及 $g$ 非负, 对几乎处处的 $t$ 有
\[
 h'(t)=g(t)y(t)\leq g(t)h(t).
\]
于是, 若令
\[
 z(t)=h(t)\exp\left(-\int_0^t g(s)\,\mathrm{d}s\right)
\]
就立即可见 $z\in W^{1,1}((0,T))$, 且
\[
 z'(t)=\bigl(h'(t)-g(t)h(t)\bigr)
 \exp\left(-\int_0^t g(s)\,\mathrm{d}s\right).
\]
因此, 对几乎处处的 $t$, 有 $z'(t)\leq0$. 由 Corollary~\ref{cor:chap2-w11-continuous-representative}, 这表明函数 $z$ 非增, 因而
\[
 z(t)\leq z(0)=h(0)=y_0,
 \qquad\forall t\in[0,T],
\]
亦即
\[
 h(t)\leq y_0\exp\left(\int_0^t g(s)\,\mathrm{d}s\right),
 \qquad\forall t\in[0,T].
\]
由于根据假设有 $y\leq h$ 几乎处处, 断言得证.
\end{Proof}

\begin{Lemma}[Uniform Gronwall lemma]
\label{lem:chap2-uniform-gronwall}
\MBUnresolvedReference{bib:source-121}{[121]}
设非负函数 $g_1,g_2\in L^1_{\mathrm{loc}}(\mathbb{R}_+)$, 并存在 $k_1,k_2$, 使
\[
 \int_t^{t+1}g_1(s)\,\mathrm{d}s\leq k_1,
 \qquad\int_t^{t+1}g_2(s)\,\mathrm{d}s\leq k_2,
 \quad\forall t\geq0.
\]
设 $y\in C^1([0,+\infty),\mathbb{R}_+)$ 且
\begin{align}
 y'(t)&\leq g_1(t)+g_2(t)y(t),\quad\text{a.e. }t\geq0,
 \label{eq:chap2-uniform-gronwall-differential}\\
 y(0)&\leq k_3,\qquad \int_t^{t+1}y(s)\,\mathrm{d}s\leq k_3,
 \quad\forall t\geq0.
 \label{eq:chap2-uniform-gronwall-average}
\end{align}
则 $y$ 在 $\mathbb{R}_+$ 上有界, 且有如下上界:
\[
 y(t)\leq(k_1+k_3)e^{k_2},\qquad\forall t\geq0.
\]
\end{Lemma}

\begin{Proof}
在 $s$ 与 $t$ 之间对式\eqref{eq:chap2-uniform-gronwall-differential}积分, 其中 $0\leq s\leq t$, 得
\[
 y(t)\leq y(s)+\int_s^t g_1(\tau)\,\mathrm{d}\tau
 +\int_s^t g_2(\tau)y(\tau)\,\mathrm{d}\tau.
\]
由 Lemma~\ref{lem:chap2-gronwall}推出
\[
 y(t)\leq\left(y(s)+\int_s^t g_1(\tau)\,\mathrm{d}\tau\right)
 \exp\left(\int_s^t g_2(\tau)\,\mathrm{d}\tau\right).
\]
\begin{itemize}
\item 当 $t\leq1$ 时, 取 $s=0$, 直接得到结论.
\item 当 $t\geq1$ 时, 取 $s\in[t-1,t]$, 并应用式\eqref{eq:chap2-uniform-gronwall-average}, 得
\[
 y(t)\leq\bigl(y(s)+k_1\bigr)e^{k_2}.
\]
保持 $t$ 固定, 对这一不等式关于 $s$ 在 $t-1$ 与 $t$ 之间积分, 得
\[
 y(t)\leq(k_3+k_1)e^{k_2}.
\]
\end{itemize}
\end{Proof}

本节最后给出一个同类型的结果, 它在研究某些非线性方程时十分有用. 这一结果类似于微分不等式之间通常的比较定理, 但其假设以更弱的积分形式表述. 这解释了为什么需要对非线性项作某种单调性假设.

\begin{Lemma}[Bihari's inequality]
\label{lem:chap2-bihari}
\MBUnresolvedReference{bib:source-16}{[16]}
设 $f:[0,+\infty)\to[0,+\infty)$ 连续且非减, 在 $(0,+\infty)$ 上为正, 并满足 $\int_1^{+\infty}1/f(x)\,\mathrm{d}x<+\infty$. 以 $F$ 表示在 $+\infty$ 处为零的 $-1/f$ 的原函数. 设 $y$ 在 $[0,+\infty)$ 上连续非负, $g\in L^1_{\mathrm{loc}}([0,+\infty))$ 非负, 且存在 $y_0>0$ 使
\[
 y(t)\leq y_0+\int_0^t g(s)\,\mathrm{d}s+\int_0^t f(y(s))\,\mathrm{d}s,
 \qquad\forall t\geq0.
\]
则方程
\begin{equation}
 T^*=F\left(y_0+\int_0^{T^*}g(s)\,\mathrm{d}s\right)
 \label{eq:chap2-bihari-critical-time}
\end{equation}
有唯一解 $T^*$, 且对每个 $T<T^*$,
\[
 \sup_{t\leq T}y(t)\leq
 F^{-1}\left(F\left(y_0+\int_0^Tg(s)\,\mathrm{d}s\right)-T\right).
\]
\end{Lemma}

\begin{Proof}
$F$ 非增且在 $+\infty$ 趋于 $0$, 而 $g$ 非负, 因而满足式\eqref{eq:chap2-bihari-critical-time}的 $T^*$ 存在且唯一. 取 $0<T<T^*$. 对所有 $t\leq T$, 由 $g$ 非负可得
\begin{equation}
 y(t)\leq y_0+\int_0^Tg(s)\,\mathrm{d}s
 +\int_0^t f(y(s))\,\mathrm{d}s.
 \label{eq:chap2-bihari-majorant}
\end{equation}
将该不等式的右端记为 $z_T(t)$. 由于 $f$ 和 $y$ 连续, 函数 $z_T$ 属于 $C^1$, 且
\[
 z_T(0)=y_0+\int_0^Tg(s)\,\mathrm{d}s.
\]
对所有 $t<T$, 有
\[
 z_T'(t)=f(y(t))\leq f(z_T(t)),
\]
因为 $f$ 非减. 注意, $z_T$ 是非减函数; 又由于 $y_0>0$, 函数 $z_T$ 不为零. 因此
\[
 \frac{z_T'(t)}{f(z_T(t))}\leq1,
 \qquad\forall t<T.
\]
在时刻 $0$ 与 $T$ 之间积分, 得
\[
 F(z_T(T))-F(z(0))\geq-T.
\]
$F$ 非增, 因而
\begin{equation}
 z_T(T)\leq F^{-1}\bigl(F(z_T(0))-T\bigr)
 =F^{-1}\left(F\left(y_0+\int_0^Tg(s)\,\mathrm{d}s\right)-T\right).
 \label{eq:chap2-bihari-bound}
\end{equation}
注意, 该式是有意义的, 因为 $T^*$ 的定义以及条件 $T<T^*$ 蕴含 $F(z_T(0))-T$ 属于 $F$ 的值域. 由式\eqref{eq:chap2-bihari-majorant}以及 $z_T$ 非减可得
\[
 y(t)\leq z_T(t)\leq z_T(T),
 \qquad\forall t<T.
\]
因此式\eqref{eq:chap2-bihari-bound}给出断言.
\end{Proof}
